% Source: https://artofproblemsolving.com/community/c3h1628055_need_a_extended_solution_to_this_problem
% Stable identity: aops-topic-1628055-q01
\textbf{(a)} Split the day at the times when someone arrives or leaves:
\begin{itemize}
\item 9:00 to 11:00: only the first group, $20$ employees;
\item 11:00 to 12:00: the first and second groups, $20+18=38$ employees;
\item 12:00 to 2:00: the second and third groups, $18+14=32$ employees.
\end{itemize}
The most employees are present from $\boxed{11{:}00\text{ AM to }12{:}00\text{ noon}}$, when $38$ are at work.

\textbf{(b)} The smallest of the three numbers is $20>8$, so at every moment of the working day more employees are present than there are computers; if everyone wants a computer, all $8$ are in use from 9:00 AM to 2:00 PM. The computers are therefore busy all day, but the demand on them (employees per computer) is greatest from 11:00 to 12:00, with $\dfrac{38}{8}=4.75$ employees per computer, compared with $2.5$ before 11:00 and $4$ after noon.
